\documentclass[UTF8,a4paper,11pt,openany]{ctexbook}
\usepackage{../common/statlearnbook}
\SLSetBookMetadata{第 5 册}{统计学习方法讲义 v2.6.0}{采样方法、主题模型与图排序}
\begin{document}
\setcounter{page}{766}
\setcounter{chapter}{32}
\setcounter{figure}{5}
\pagestyle{headings}
\chaptermark{马尔可夫链蒙特卡罗法}
\sectionmark{接受率推导}
设比值
\[
r=\frac{\pi(y)q(y,x)}{\pi(x)q(x,y)}.
\]
MH 接受概率为 $\alpha(x,y)=\min\{1,r\}$。图\ref{fig:V5-C03-acceptance-function}同时画出比值小于与大于 $1$ 的两段；“截到 $1$”来自概率上界，而不是数值补丁。
% v2.3.1 figure UID: FIG-P603-01
% v2.6.0 reverse redraw for FIG-P603-01: explicit hierarchy and collision-free labels.
\tikzset{slfig-FIG-P603-01/.style={font=\fontsize{9.2pt}{11.0pt}\selectfont,every path/.append style={line cap=round,line join=round}}}
\pgfplotsset{slfig-FIG-P603-01-axis/.style={
  tick label style={font=\fontsize{8.5pt}{10.0pt}\selectfont},
  label style={font=\fontsize{9.2pt}{11.0pt}\selectfont},
  axis line style={line width=.65pt},tick style={line width=.55pt},clip=false}}
\begin{figure}[htbp]
\centering
\begin{tikzpicture}[slfig-FIG-P603-01,every node/.style={font=\fontsize{9.2pt}{11.0pt}\selectfont,align=center}]
\begin{axis}[slfig-FIG-P603-01-axis,
  name=af,width=8.6cm,height=5.7cm,
  xmin=0,xmax=3,ymin=0,ymax=1.15,axis lines=left,
  xlabel={$r$},
  ylabel={$\alpha(x,y)$},
  xtick={0,1,2,3},ytick={0,.5,1},clip=false
]
  \addplot[draw=none,fill=SLBlue!7] coordinates {(0,0)(1,1)(3,1)(3,0)} \closedcycle;
  \addplot[draw=SLBlue,line width=1.02pt,domain=0:1,samples=181]{x};
  \addplot[draw=SLBlue,line width=1.02pt,domain=1:3,samples=181]{1};
  \draw[draw=SLGray,line width=.55pt,densely dashed] (axis cs:1,0)--(axis cs:1,1)--(axis cs:0,1);
  \addplot[only marks,mark=*,mark size=2.4pt,draw=SLGold,fill=SLGold] coordinates {(1,1)};
  \node[align=center,text=SLBlue] at (axis cs:.68,.27) {$r<1$：\\按比例\\接受};
  \node[text=SLBlue] at (axis cs:2.08,.80) {$r\ge1$：必然接受};
  \node[anchor=south,inner sep=0pt,yshift=2.5mm,text=SLGold] at (axis cs:1,1) {折点};
\end{axis}
\node[right=8mm of af.east,draw=SLRuleGray,rounded corners=2pt,fill=white,line width=.58pt,
  font=\fontsize{9.2pt}{11.0pt}\selectfont,align=left,inner xsep=3mm,inner ysep=2.3mm]
  {$\displaystyle r=\frac{\pi(y)q(y,x)}{\pi(x)q(x,y)}$\\[3pt]
   独立提议：$\displaystyle r=\frac{w(y)}{w(x)}$};
\end{tikzpicture}
\caption{MH接受概率是比值 $r$ 的截断函数 $\alpha=\min\{1,r\}$，其中 $r=\pi(y)q(y,x)/[\pi(x)q(x,y)]$；独立提议时可写为 $r=w(y)/w(x)$}\label{fig:V5-C03-acceptance-function}
\end{figure}

\paragraph{常见特例。}
若 $q(x,y)=q(y,x)$，提议项消去；若 $q(x,y)=q(y)$ 与当前状态无关，则比值可写为独立提议权重之比 $w(y)/w(x)$。
\end{document}
